\documentclass[12pt, letterpaper]{article}
\usepackage[margin=1.2in]{geometry}
\usepackage{ebgaramond}
\usepackage[T1]{fontenc}
\usepackage{microtype}
\usepackage{setspace}
\usepackage{parskip}
\usepackage{titling}
\usepackage{fancyhdr}

\setlength{\parindent}{2em}
\onehalfspacing

\pagestyle{fancy}
\fancyhf{}
\fancyhead[L]{\small\textit{The Canopy}}
\fancyhead[R]{\small\textit{NASA Mission Series v1.5}}
\fancyfoot[C]{\small\thepage}
\renewcommand{\headrulewidth}{0.4pt}

\title{\Huge\textbf{The Canopy}\\[0.5em]\large\textit{A short story from Mission 1.5 --- Pioneer 4}}
\author{NASA Mission Series\\[0.3em]\small Miles Tiberius Foxglove --- Mukilteo, WA}
\date{March 29, 2026}

\begin{document}
\maketitle
\thispagestyle{empty}

\vspace{1em}
\noindent\rule{\textwidth}{0.4pt}
\vspace{1em}

\noindent I have been here for nine hundred and seventy-three years. I know this because I have counted every one of them --- not in the way you count, with numerals and calendars, but in the way a river counts its stones: by laying them down, one after another, each one a ring of wood laid beneath my bark in the long light of summer and sealed shut by the frosts of autumn. Nine hundred and seventy-three rings. You could read me like a book, if you knew how to look.

I began as a seed no larger than your thumbnail, with a single papery wing. My mother was ninety meters tall the year she released me --- one of a hundred thousand seeds she dropped that October, spiraling out of cones that the crossbills had been working for weeks. The wind took me. I did not choose my direction. The wing caught an updraft along the ridgeline and carried me forty meters from her crown, past the edge of her shadow, past the understory where the hemlocks were already tall and patient. I landed on a nurse log --- a Douglas-fir that had fallen eighty years earlier and was halfway through its long return to soil. The wood was soft, brown, furred with moss, and full of water. I settled into a crack in the bark.

The first year was the hardest. I say that, but every year since has had its own claim on hardship, and I suspect the very last year --- whenever it arrives --- will contest the title. But the first year was hard because the forest floor receives two percent of the sunlight that reaches the canopy. Two percent. Stand where I stand now and look up: you will see a green ceiling a hundred meters above you, and through it, coins of light that shift with the wind and vanish with the clouds. Two percent of the sun is all that reaches the ground, and I needed light to photosynthesize, and I needed to photosynthesize to grow, and I needed to grow to reach the light. The problem was circular, and the circle was a trap that killed nine hundred and ninety-nine of my thousand siblings in their first decade.

I survived because the nurse log gave me a head start --- a meter of elevation above the fern layer, a substrate that held moisture through the summer droughts, a position directly beneath a canopy gap where a hemlock had fallen two winters before my birth. That gap let in nine percent of the light instead of two. Nine percent does not sound like much, but it was the difference between slow death and slow growth. I grew twelve centimeters my first year. My siblings on the flat ground grew three.

\begin{center}
$\bullet\quad\bullet\quad\bullet$
\end{center}

There was a time, long before my seed left my mother's cone, when five small objects were launched from a place called Cape Canaveral toward the sky above Earth. I know about them the way I know about the salmon in the rivers two valleys over --- not through direct experience, but through the network. The mycorrhizal threads that link my roots to every other tree in this stand carry more than nutrients. They carry signals. Chemical whispers. And in those whispers, over decades, the forest assembles a picture of the world beyond its canopy. We are slow, but we are thorough.

The first object exploded seventy-seven seconds after it left the ground. I understand explosions. I have been struck by lightning twice --- once in my three hundred and twelfth year and once in my six hundred and forty-first. Both times the bolt ran down my trunk and blew a strip of bark from crown to root. Both times I survived because my sapwood is wet and my bark is thick. The rocket had no bark.

The second object rose to a height of one hundred and thirteen thousand kilometers. One hundred and thirteen thousand. I have spent nine centuries reaching one hundred meters. The rocket reached a hundred thousand kilometers in hours. But then it stopped climbing and fell back, because it did not have enough speed to leave. Not enough speed. I understand that too. In my forty-third year, a windstorm blew down the hemlock that had been shading my canopy gap. The gap opened wider. The light increased from nine percent to fourteen percent. I grew a meter and a half that year --- the fastest I have ever grown. I was racing for the canopy the way the rocket was racing for escape, and for exactly the same reason: below the threshold, you fall back. Above it, you keep going.

The third object barely left the ground. Its upper stage --- a thing like a secondary trunk that was supposed to carry it higher --- failed to ignite. It reached one thousand five hundred and fifty kilometers. A regression. I have had regression years too. In my hundred and nineteenth year, a root rot fungus --- \textit{Phellinus weirii}, the Douglas-fir's quiet nemesis --- killed a third of my root system. My crown thinned. My growth rings that year are barely visible: a hair's width of wood laid down by a tree fighting for its life. I survived because the mycorrhizal network rerouted nutrients through neighboring trees, feeding me until I could regrow the roots. The rocket had no network.

The fourth object came within seventy-three meters per second of escape. Seventy-three meters per second, out of eleven thousand needed. Its fuel ran out three and a half seconds too early. It rose to a hundred and two thousand kilometers, found the second radiation belt that the second object had only hinted at, and fell back. I understand seventy-three meters per second. In the currency of my species, seventy-three meters per second is the difference between a growth ring of four millimeters (good year, ample rain, canopy gap open) and a growth ring of one millimeter (drought year, gap closing). Compounded over decades, those three millimeters per year determine whether you break through the canopy or die in the shade. The rocket needed three more seconds of fuel. I needed three more millimeters per year.

\begin{center}
$\bullet\quad\bullet\quad\bullet$
\end{center}

The fifth object launched on March 3, 1959. By then I was already nine hundred and six years old, and my trunk was nearly four meters in diameter. I had broken through the canopy six hundred years earlier, back when the concept of ``escape velocity'' was not yet a phrase in any language, back when the forest I live in was being shaped by fire and flood and the slow patient work of centuries. I had reached the light. I had been self-sustaining for six hundred years. The sun hit my crown directly, and I photosynthesized at the maximum rate my needles could sustain, and I grew outward instead of upward, thickening, deepening, becoming the architecture that everything else in this forest depends on.

The fifth object reached escape velocity.

It did not come back.

I know what that means, because I have never come back either. Once I broke through the canopy, the understory became a place I send resources to, not a place I live in. My lower branches died long ago --- they were shaded out by my own crown, and I let them go. My energy goes upward and outward: to new needles in the crown, to seeds in the cones, to the mycorrhizal network through my roots. I am not an understory tree anymore. I escaped, and the escape was permanent. The light up here is not the thin filtered light of the forest floor. It is the full light of the sun, and it has been sustaining me for six centuries, the way the sun has been sustaining that small metal cone for six decades.

They named the object Pioneer~4. They named me \textit{Pseudotsuga menziesii}, which means ``false hemlock, of Menzies,'' after a man named Archibald Menzies who saw my kind on Vancouver Island in 1791 and decided I needed a European name. I did not. The Snohomish people who lived beside me for thousands of years before Menzies arrived had their own name for me. But names are like seeds: they land where the wind drops them, and the one that takes root is the one you are stuck with. I am a false hemlock named after a Scottish surgeon, and I am the largest species of tree in the Pacific Northwest, and I have been here for nearly a millennium, and I do not need a name to be what I am.

\begin{center}
$\bullet\quad\bullet\quad\bullet$
\end{center}

I want to tell you about what lives on me. Not in me --- though there is a red tree vole in my crown that has not touched the ground in three generations, and a family of Red-breasted Nuthatches in the cavity where my six-hundredth-year branch broke off. On me. My bark is a landscape. Thirty centimeters thick in places, deeply furrowed, ridged like a mountain range seen from above. In those furrows live mosses, lichens, liverworts, and insects that have never been to the ground. One lichen in particular --- \textit{Usnea longissima}, Old Man's Beard, pale green-gold filaments hanging from my mid-canopy branches like hair --- has been growing on me for longer than Pioneer~4 has been orbiting the Sun. The longest strands are a hundred years old, growing at two millimeters a year, a meter of patience. They are the organism of the previous mission --- v1.4, the lichen that is two organisms in one.

The Nuthatch comes down my trunk headfirst. It is the only bird that does this --- walks straight down the bark, peering into the crevices that upward-climbing birds never see. It finds insects that no other bird finds, in positions that no other bird reaches. It sees me from an angle that even I cannot see myself. When it calls --- \textit{yank, yank, yank} --- the sound is thin and nasal and it cuts through the forest like a tin horn. It does not carry far, but it carries far enough. A hundred meters, maybe two hundred, before the forest absorbs it. The Nuthatch does not need to be heard at six hundred thousand kilometers. It needs to be heard by the next Nuthatch, on the next Douglas-fir, a hundred meters away. Scale your transmitter to your audience.

I am an ecosystem. That is not a metaphor. I am literally an ecosystem --- a platform on which hundreds of other species live, feed, reproduce, and die. My trunk is the structure. My canopy is the energy source. My roots are the communication network. My bark is the terrain. My fallen branches are the nutrient cycle. When I eventually fall --- and I will; nothing made of cellulose lasts forever, not even cellulose four meters in diameter --- I will become a nurse log, the way my mother's mother became the nurse log that I sprouted on. I will spend the next four hundred years returning to soil, and during those four hundred years I will support more biomass as a dead log than most living trees support in their entire lives. The fall is not the end. The fall is the transition from vertical ecosystem to horizontal ecosystem.

\begin{center}
$\bullet\quad\bullet\quad\bullet$
\end{center}

There is a man born on the same day that Pioneer~4 launched, but a hundred and twelve years earlier. March 3, 1847. Alexander Graham Bell. He invented a device that carried the human voice through a copper wire, across distances that the voice alone could not reach. The voice is a vibration in air. The wire carries a vibration in current. The principle is the same as the mycorrhizal network that connects me to the hemlock forty meters to my west: encode the signal in a medium that can carry it farther than the original medium allows. Bell sent voice through wire. Pioneer~4 sent data through vacuum. I send carbon through fungi. The Nuthatch sends territory through air. All of us are solving the same problem: how do you reach across a distance that your native signal cannot cross?

Bell would have understood the Nuthatch, I think. He spent his life studying sound --- how it is produced, how it travels, how it is received. His father developed Visible Speech, a system for transcribing any sound the human voice could make. His mother was deaf, and he married a woman who was deaf, and his entire career was shaped by the problem of bridging the gap between a person who speaks and a person who cannot hear. The telephone was not a gadget. It was a bridge. Pioneer~4's transmitter was a bridge. My mycorrhizal network is a bridge. The Nuthatch's call is a bridge. We are all building bridges across distances that silence would otherwise own.

\begin{center}
$\bullet\quad\bullet\quad\bullet$
\end{center}

I want to tell you one more thing, and then I will be quiet, because I am a tree and quiet is what I do best.

The canopy is not the destination.

I spent my first three hundred years reaching it. I have spent the last six hundred years living in it. But the canopy is not where I was trying to go. The canopy is where the sky begins. Above me there is nothing but air, and above the air there is nothing but space, and in that space there is a six-kilogram cone of aluminum that left Earth sixty-seven years ago and has been orbiting the Sun ever since, silent, tumbling, permanent. It broke through the gravity well the way I broke through the canopy: not to stop there, but to keep going.

The canopy is not a ceiling. It is a floor. The floor of the sky. When I reached it, I did not stop growing. I grew wider. I grew stronger. I grew more branches, supported more species, sent more carbon into the network, dropped more seeds for the next generation. The canopy was not the answer. The canopy was the question: now that you are here, what will you do with the light?

Pioneer~4 answered that question with 82 hours of data transmitted across 655,300 kilometers before its batteries died and it went silent forever. I have answered that question with nine hundred and seventy-three years of growth, a trunk that four people could not encircle with their arms, a crown that shelters a quarter-acre of forest floor, and a root system that feeds the network that feeds the forest that feeds the rivers that feeds the ocean that feeds the rain that falls on me.

The canopy is not the destination. It is where the sky begins.

Come stand at my base. Look up. The trunk narrows for a hundred meters above you, furrowed bark receding into green light, and at the top, so far away that the branches are just dark lines against brightness, you can see the sky. That is where the spacecraft went. That is where I am. That is where every seed that catches the right wind and every rocket that burns long enough eventually arrives: the place where the canopy ends and the light is yours.

\vspace{2em}
\begin{center}
\rule{3cm}{0.4pt}
\end{center}
\vspace{1em}

\begin{center}
\small\textit{NASA Mission Series v1.5 --- Pioneer 4 + Douglas-fir}\\
\small Miles Tiberius Foxglove --- Mukilteo, WA --- March 29, 2026
\end{center}

\end{document}
